% Part: sets-functions-relations
% Chapter: sets
% Section: pairs-and-products

\documentclass[../../../include/open-logic-section]{subfiles}

\begin{document}

\olfileid[pt]{sfr}{set}{pai}
\olsection{Pares, Tuplas, Produtos Cartesianos}

\begin{explain}
Da extensionalidade segue-se que os conjuntos não têm ordem entre os seus
elementos. Por isso, se queremos representar ordem, usamos \emph{pares ordenados}
$\tuple{x, y}$. Num par não ordenado $\{x, y\}$, a ordem não
importa: $\{x, y\} = \{y, x\}$. Num par ordenado, importa: se $x
\neq y$, então $\tuple{x, y} \neq \tuple{y, x}$.

Como devemos pensar os pares ordenados na teoria dos conjuntos? O essencial é
preservar a ideia de que pares ordenados são idênticos se e só se
partilham o mesmo primeiro elemento e o mesmo segundo elemento, O.s.:
\[
  \tuple{a, b}= \tuple{c, d}\text{ se e só se }a = c \text{ e }b=d.
\]
Podemos definir pares ordenados na teoria dos conjuntos usando a definição de
Wiener--Kuratowski.
\end{explain}

\begin{defn}[Par ordenado]\ollabel{wienerkuratowski}
	$\tuple{a, b} = \{\{a\}, \{a, b\}\}$.
\end{defn}

\begin{prob}
	Usando \olref[sfr][set][pai]{wienerkuratowski}, prove que $\tuple{a,
	b}= \tuple{c, d}$ se e só se $a = c$ e $b=d$.
\end{prob}

\begin{explain}
Tendo fixada uma definição de par ordenado, podemos usá-la para definir
outros conjuntos. Por exemplo, por vezes queremos sequências ordenadas de
mais de dois objetos, p.ex., \emph{triplos} $\tuple{x, y, z}$,
\emph{quádruplos} $\tuple{x, y, z, u}$, e assim por diante.  Podemos pensar nos
triplos como pares ordenados especiais, em que o primeiro elemento é ele próprio um
par ordenado: $\tuple{x, y, z}$ é $\tuple{\tuple{x, y},z}$. O mesmo
vale para quádruplos: $\tuple{x,y,z,u}$ é
$\tuple{\tuple{\tuple{x,y},z},u}$, e assim por diante. Em geral, falamos de
\emph{$n$-tuplas ordenadas} $\tuple{x_1, \dots, x_n}$.

Certos conjuntos de pares ordenados, ou outras $n$-tuplas ordenadas, serão úteis.
\end{explain}

\begin{defn}[Produto cartesiano]
Dados conjuntos $A$ e $B$, o seu \emph{produto cartesiano} $A \times B$ é
definido por
\[
  A \times B = \Setabs{\tuple{x, y}}{x \in A \text{ e } y \in B}.
\]
\end{defn}

\begin{ex}
Se $A = \{0, 1\}$, e $B = \{1, a, b\}$, então o seu produto é
\[
A \times B = \{ \tuple{0, 1}, \tuple{0, a}, \tuple{0, b},
    \tuple{1, 1}, \tuple{1, a}, \tuple{1, b} \}.
\]
\end{ex}

\begin{ex}
Se $A$ é um conjunto, o produto de $A$ consigo próprio, $A \times A$, também se
escreve~$A^2$. É o conjunto de \emph{todos} os pares $\tuple{x, y}$ com
$x, y \in A$. O conjunto de todos os triplos $\tuple{x, y, z}$ é $A^3$, e
assim por diante. Podemos dar uma definição recursiva:
\begin{align*}
  A^1 & = A\\
  A^{k+1} & = A^k \times A
\end{align*}
\end{ex}

\begin{prob}
Liste todos os !!{element}s de $\{1, 2, 3\}^3$.
\end{prob}

\begin{prop}\ollabel{cardnmprod}
Se $A$ tem $n$ !!{element}s e $B$ tem $m$ !!{element}s, então $A
\times B$ tem $n\cdot m$ elementos.
\end{prop}

\begin{proof}
Para todo !!{element}~$x$ em~$A$, há $m$ !!{element}s da
forma $\tuple{x, y} \in A \times B$. Seja $B_x = \Setabs{\tuple{x, y}}{y
  \in B}$. Como sempre que $x_1 \neq x_2$, $\tuple{x_1, y} \neq
\tuple{x_2, y}$, tem-se $B_{x_1} \cap B_{x_2} = \emptyset$. Mas se $A = \{x_1,
\dots, x_n\}$, então $A \times B = B_{x_1} \cup \dots \cup B_{x_n}$, e portanto tem
$n\cdot m$ !!{element}s.

Para visualizar isto, dispomos os !!{element}s de~$A \times B$ numa grade:
\[
\begin{array}{rcccc}
  B_{x_1} = & \{\tuple{x_1, y_1} & \tuple{x_1, y_2} & \dots & \tuple{x_1, y_m}\}\\
  B_{x_2} = & \{\tuple{x_2, y_1} & \tuple{x_2, y_2} & \dots & \tuple{x_2, y_m}\}\\
  \vdots & & \vdots\\
  B_{x_n} = & \{\tuple{x_n, y_1} & \tuple{x_n, y_2} & \dots & \tuple{x_n, y_m}\}
\end{array}
\]
Como os $x_i$ são todos distintos, e os $y_j$ são todos distintos, quaisquer
dois pares nesta grade são distintos, e há $n\cdot m$
deles.
\end{proof}

\begin{prob}
Mostre, por indução em~$k$, que para todo $k \ge 1$, se $A$ tem $n$
!!{element}s, então $A^k$ tem $n^k$ !!{element}s.
\end{prob}

\begin{ex}
Se $A$ é um conjunto, uma \emph{palavra} sobre~$A$ é qualquer sequência de
!!{element}s de~$A$. Uma sequência pode ser vista como uma $n$-tupla de
!!{element}s de~$A$. Por exemplo, se $A = \{a, b, c\}$, então a
sequência ``$bac$'' pode ser vista como o triplo~$\tuple{b, a, c}$.
Palavras, O.s.\@, sequências de símbolos, são de importância crucial na
ciência da computação. Por convenção, contamos os !!{element}s de~$A$ como
sequências de comprimento~$1$, e $\emptyset$ como a sequência de comprimento~$0$.
O conjunto de \emph{todas} as palavras sobre~$A$ é então
\[
A^* = \{\emptyset\} \cup A \cup A^2 \cup A^3 \cup \dots
\]
\end{ex}

\end{document}
